\documentclass{article}

\usepackage{booktabs}
\usepackage{tabularx}
\usepackage{hyperref}

\hypersetup{
    colorlinks=true,       % false: boxed links; true: colored links
    linkcolor=red,          % color of internal links (change box color with linkbordercolor)
    citecolor=green,        % color of links to bibliography
    filecolor=magenta,      % color of file links
    urlcolor=cyan           % color of external links
}

\title{Hazard Analysis\\\progname}

\author{\authname}

\date{}

\input{../Comments.text}
\input{../Common.text}

\begin{document}

\maketitle
\thispagestyle{empty}

~\newpage

\pagenumbering{roman}

\begin{table}[hp]
\caption{Revision History} \label{TblRevisionHistory}
\begin{tabularx}{\textwidth}{llX}
\toprule
\textbf{Date} & \textbf{Developer(s)} & \textbf{Change}\\
\midrule
Sept 30 & Sankalpa Chhetri Dhakal & Completed Introduction\\
Oct 6 & Sankalpa Chhetri Dhakal & Completed Scope and Purpose and System Boundaries\\
... & ... & ...\\
\bottomrule
\end{tabularx}
\end{table}

~\newpage

\tableofcontents

~\newpage

\pagenumbering{arabic}

\wss{You are free to modify this template.}

\section{Introduction}

\wss{You can include your definition of what a hazard is here.}

\section{Scope and Purpose of Hazard Analysis}

\section{System Boundaries and Components}

The system boundary includes the software and hardware directly involved in
running and interacting with the FPV drone training simulator. This includes
the VR headset, input devices, host computer, Unity application, drone
simulation, training functionality, user interface, and locally stored data.
The user and their physical surroundings interact with the system and are
considered when identifying hazards, but they are not components of the
software system itself. Real FPV drones and real-world drone operation are
outside the system boundary.

The system consists of several major software and hardware components. These
components are considered independently during hazard identification because a
failure in one component may affect the safety or correctness of the overall
system.

\subsection{VR Headset}

The VR headset, expected to be a \textbf{Meta Quest 3S}, provides the user's
visual representation of the simulated environment and tracks the position and
orientation of the user's head.

Potential hazards associated with this component include loss of tracking,
visual latency, incorrect orientation, poor headset fit, and the user becoming
unaware of their physical surroundings.

\subsection{VR Controllers / Input Devices}

The input subsystem allows the user to control the simulated drone. Depending
on the final implementation, input may be provided through VR controllers, a
game controller, or an external FPV-style radio controller.

Failures may include disconnected controllers, incorrect input mappings, stuck
inputs, input delay, or incorrect calibration.

\subsection{Host Computer}

A computer executes the Unity application and associated simulation software.

Failures may include insufficient hardware performance, overheating, loss of
power, driver issues, or resource exhaustion.

\subsection{Unity Application}

Unity provides the primary execution environment for the game.

This component is responsible for scene management, game logic, user
interaction, rendering, and communication between major software modules.

\subsection{Drone Physics Simulation}

The drone physics subsystem models properties such as:

\begin{itemize}
    \item thrust;
    \item gravity;
    \item acceleration;
    \item angular velocity;
    \item momentum;
    \item collisions; and
    \item drone response to pilot input.
\end{itemize}

Incorrect behaviour in this component represents an important hazard because
users are expected to learn flying behaviour from the simulator.

\subsection{VR Rendering and XR Subsystem}

This subsystem renders the virtual environment and communicates with the VR
headset.

Failures may result in visual artifacts, low frame rate, latency, incorrect
camera movement, or loss of tracking.

\subsection{Training and Scoring System}

The training subsystem may contain tutorials, objectives, checkpoints, scoring
metrics, crash detection, and user performance feedback.

Incorrect scoring or feedback could indicate that unsafe or incorrect flying
behaviour is acceptable.

\subsection{User Interface}

The interface includes menus, settings, warnings, calibration options, pause
controls, and training instructions.

Poor interface design could prevent the user from understanding warnings or
exiting the simulation quickly.

\subsection{Data Storage}

The system may locally store information such as:

\begin{itemize}
    \item user settings;
    \item controller configuration;
    \item training progress;
    \item scores; and
    \item application configuration.
\end{itemize}

Data corruption or unauthorized modification could affect system operation or
training results.

\section{System Boundaries and Components}

\wss{Dividing the system into components will help you brainstorm the hazards.
You shouldn't do a full design of the components, just get a feel for the major
ones.  For projects that involve hardware, the components will typically include
each individual piece of hardware.  If your software will have a database, or an
important library, these are also potential components.}

\section{Critical Assumptions}

\wss{These assumptions that are made about the software or system.  You should
minimize the number of assumptions that remove potential hazards.  For instance,
you could assume a part will never fail, but it is generally better to include
this potential failure mode.}

\section{Failure Mode and Effect Analysis}

\wss{Include your FMEA table here. This is the most important part of this document.}
\wss{The safety requirements in the table do not have to have the prefix SR.
The most important thing is to show traceability to your SRS. You might trace to
requirements you have already written, or you might need to add new
requirements.}
\wss{If no safety requirement can be devised, other mitigation strategies can be
entered in the table, including strategies involving providing additional
documentation, and/or test cases.}

\section{Safety and Security Requirements}

\wss{Newly discovered requirements.  These should also be added to the SRS.  (A
rationale design process how and why to fake it.)}

\section{Roadmap}

\wss{Which safety requirements will be implemented as part of the capstone timeline?
Which requirements will be implemented in the future?}

\newpage{}

\section{Appendix --- Generative AI Use Disclosure}

\wss{ChatGPT (version 5) was used to brainstorm ideas for the template for this disclose section.}

\subsection{AI Tools Used}

\wss{Give the name of the AI tool(s) used and the version.}

\subsection{How and Where Used}

\wss{Explain what the tool(s) were used for.  Examples include: brainstorming,
explaining a technical concept, reviewing the document, generating or modifying
text, generating or modifying diagrams, identifying errors or omissions or
inconsistencies.}

\wss{For each of the uses, identify which parts of the document were affected.}

\wss{You do not need to report every interaction, but you should disclose the 
uses that materially contributed to your work}

\subsection{Verification of AI-Generated Content}

\wss{\begin{itemize}
    \item Describe how the team checked AI-generated or AI-assisted material
    before including it in the document.
    \item In particular, verify factual claims, technical assertions,
    requirements, calculations, references, and other information for which an
    AI system may produce incorrect or fabricated results.
    \item \textbf{The team is responsible for the accuracy, completeness,
consistency, and appropriateness of the submitted document, regardless of
whether material was generated by a human or an AI system.}
\end{itemize}}

\newpage{}

\section{Appendix --- Reflection}

\wss{Not required for CAS 741}

\input{../Reflection.text}

\begin{enumerate}
    \item What went well while writing this deliverable? 
    \item What pain points did you experience during this deliverable, and how
    did you resolve them?
    \item Which of your listed risks had your team thought of before this
    deliverable, and which did you think of while doing this deliverable? For
    the latter ones (ones you thought of while doing the Hazard Analysis), how
    did they come about?
    \item Other than the risk of physical harm (some projects may not have any
    appreciable risks of this form), list at least 2 other types of risk in
    software products. Why are they important to consider?
\end{enumerate}

\end{document}